

\documentclass[../../main.tex]{subfiles}




\begin{document}

\section{Analytical phonon integration in Hubbard--Holstein VMC}

\paragraph{Abstract}
A conditional phonon state can be contracted analytically before sampling the
fermionic configurations. For coherent packets, normally ordered matrix
elements follow from the coherent-state eigenvalue equation. For squeezed
packets, they follow by differentiating one overlap kernel. The same kernels
give the energy gradient, the quantum geometric tensor used in stochastic
reconfiguration, and the fermion--phonon R\'enyi entropy.


\subsection{Conditional state and sampling}
\label{sec:state}

Let \(s=(n_{i\sigma})\) be a fermionic Fock configuration and let
\(f(s)\) be its Slater, backflow, or other variational amplitude. For
\(N\) phonon modes and \(K\) packets, define
\begin{align}
 \ket{\phi_s} &= \sum_{k=1}^K W_k(s)\ket{G_k(s)}, &
 \ket{G_k(s)} &= \bigotimes_{i=1}^{N}\ket{G_{i,k}(s)}, \label{eq:packets}\\
 Z_s &= \braket{\phi_s}{\phi_s}, &
 \ket{\chi_s} &= Z_s^{-1/2}\ket{\phi_s}. \label{eq:normalized-conditional}
\end{align}
Each \(\ket{G_{i,k}}\) is normalized. The complex weights describe a
superposition, so pairwise interference contributes to \(Z_s\).
The product over sites is a tensor product within one packet. The
superposition can correlate different phonon modes.

The normalized conditional states give the following expression for the
trial state and its sampling measure.
\begin{equation}
 \ket{\Psi}=\sum_s A(s)\ket{s}\otimes\ket{\chi_s},
 \qquad
 \mathcal{N}=\langle\Psi\mid\Psi\rangle=\sum_s |A(s)|^2,\qquad
 p(s)=|A(s)|^2/\mathcal{N}.
 \label{eq:state}
\end{equation}
The notes use two conventions for \(f\).
\begin{description}
 \item[Normalized conditional ansatz.] If
 \(\ket{\Psi}=\sum_s f(s)\ket{s}\otimes\ket{\chi_s}\), then \(A(s)=f(s)\)
 and the sampler uses \(p(s)\propto |f(s)|^2\).
 \item[Unnormalized conditional ansatz.] If
 \(\ket{\Psi}=\sum_s f(s)\ket{s}\otimes\ket{\phi_s}\), then
 \(A(s)=f(s)\sqrt{Z_s}\) and the sampler uses \(p(s)\propto |f(s)|^2Z_s\).
\end{description}
These describe the same state after the corresponding redefinition of the
fermionic amplitude. For a fixed \(f\), they generally describe different
states. For a symmetric Metropolis proposal the acceptance ratio before
clipping at one is \(|A(s')/A(s)|^2\). The formulas using division by \(A\)
or \(Z_s\) are evaluated where the corresponding quantities are nonzero.
A zero unnormalized component can be kept as a zero vector in a connected
ket without constructing its normalized state.

\subsection{Phonon overlap and operator kernels}
\label{sec:kernels}

For a phonon operator \(B\), introduce the matrix and scalar kernels
\begin{align}
 [G_B(s,s')]_{\ell k} &= \mel{G_\ell(s)}{B}{G_k(s')},\\
 K_B(s,s') &= W(s)^\dagger G_B(s,s')W(s'), &
 C_B(s,s') &= \frac{K_B(s,s')}{\sqrt{Z_s Z_{s'}}}.
 \label{eq:kernels}
\end{align}
For \(B=\id\), write \(G\), \(K_{\id}\), and \(C\). In particular,
\(Z_s=K_{\id}(s,s)\) and \(C(s,s')=\braket{\chi_s}{\chi_{s'}}\).
The same-configuration matrix \(G(s,s)\) is a positive-semidefinite
Hermitian Gram matrix with unit diagonal. The cross-configuration matrix
satisfies \(G(s',s)=G(s,s')^\dagger\). For an operator acting on one site,
the tensor product gives
\begin{equation}
 [G_{B_i}(s,s')]_{\ell k}
 =\mel{G_{i,\ell}(s)}{B_i}{G_{i,k}(s')}
   \prod_{j\ne i}\braket{G_{j,\ell}(s)}{G_{j,k}(s')}.
 \label{eq:site-kernel}
\end{equation}

\subsubsection{Coherent packets}
\label{sec:coh}

For one mode, write
\(\ket{\alpha}=D(\alpha)\ket{0}\), with
\(D(\alpha)=\exp(\alpha b^\dagger-\alpha^* b)\). The scalar commutator in
the Baker--Campbell--Hausdorff identity gives
\(\ket{\alpha}=e^{-|\alpha|^2/2}e^{\alpha b^\dagger}\ket 0\).
Contracting the two exponentials gives the overlap. Applying
\(b\ket{\alpha}=\alpha\ket{\alpha}\) and its adjoint gives every normally
ordered matrix element \cite{glauber1963}. For bra \(L\) and ket \(R\),
\begin{align}
 h_{LR} &= \braket{\alpha_L}{\alpha_R}
     = \exp\!\left(-\frac{|\alpha_L|^2+|\alpha_R|^2}{2}
                    +\alpha_L^*\alpha_R\right), \label{eq:coh-overlap}\\
 \mel{\alpha_L}{(b^\dagger)^m b^n}{\alpha_R}
     &= h_{LR}(\alpha_L^*)^m\alpha_R^n. \label{eq:coh-monomial}
\end{align}
In particular, the factors multiplying \(h_{LR}\) for \(b^\dagger b\)
and \(b+b^\dagger\) are \(\alpha_L^*\alpha_R\) and
\(\alpha_R+\alpha_L^*\). For real displacements
\(h_{LR}=\exp[-(\alpha_L-\alpha_R)^2/2]\).
With complex displacements its phase is
\(\Ima(\alpha_L^*\alpha_R)\).

\subsubsection{Displaced squeezed packets}
\label{sec:sq}

Let \(\zeta=r e^{\ii\vartheta}\) be the unitary squeezing parameter and
\(\alpha\) the physical displacement. The linear and quadratic
coefficients in the creation-operator exponential are \(\beta\) and \(q\).
\begin{align}
 S(\zeta)&=\exp\!\left[\tfrac12(\zeta^* b^2-\zeta(b^\dagger)^2)\right],
 \qquad q=-e^{\ii\vartheta}\tanh r, \label{eq:squeeze}\\
 \ket{G}&=D(\alpha)S(\zeta)\ket{0}
      =c\,\exp\!\left[\beta b^\dagger+\tfrac12 q(b^\dagger)^2\right]\ket{0},\\
 \beta&=\alpha-q\alpha^*, \qquad
 c=\frac{\exp[-|\alpha|^2/2+q(\alpha^*)^2/2]}{\sqrt{\cosh r}}.
 \label{eq:gaussian-state}
\end{align}
To obtain this form, the Bogoliubov transform gives
\(S b S^\dagger=(b-qb^\dagger)\cosh r\), so
\((b-qb^\dagger)S\ket{0}=0\). Disentangling \(S\) fixes the prefactor of
the squeezed vacuum. The displacement then shifts \(b^\dagger\) to
\(b^\dagger-\alpha^*\) and maps the vacuum to \(\ket{\alpha}\).
\Cref{app:squeeze} gives the operator identities explicitly. In particular,
\(b\ket{G}=(\beta+q b^\dagger)\ket{G}\).

Put \(x=\beta_L^*\), \(y=\beta_R\), and \(d=1-q_L^*q_R\).
The single-mode overlap has the form
\begin{equation}
 h_{LR}=c_L^*c_R\,\mathcal{F}(x,y),\qquad
 \mathcal{F}(x,y)=\frac{1}{\sqrt{d}}\,
 \exp\!\left[\frac{xy+\tfrac12 q_R x^2+\tfrac12 q_L^*y^2}{d}\right].
 \label{eq:gauss-overlap}
\end{equation}
The square root is continued from \(d=1\). For finite squeezing
\(|q_{L,R}|<1\), so \(d\) lies in the right half-plane. One short way to get
\(\mathcal{F}\) is to differentiate the overlap of the bare exponentials
in \cref{eq:gaussian-state}, keeping \(q_L^*,q_R\) fixed and treating \(x,y\)
as independent. The equation \(b\ket{G}=(\beta+qb^\dagger)\ket{G}\) and
its bra counterpart give
\(\partial_x\mathcal{F}=(y+q_R\partial_y)\mathcal{F}\) and
\(\partial_y\mathcal{F}=(x+q_L^*\partial_x)\mathcal{F}\).
Solving these two equations gives
\begin{equation}
 u\equiv\partial_x\log\mathcal{F}=\frac{y+q_Rx}{d},\qquad
 v\equiv\partial_y\log\mathcal{F}=\frac{x+q_L^*y}{d}.
 \label{eq:uv}
\end{equation}
Integrating gives the exponent in \cref{eq:gauss-overlap}.
At \(x=y=0\), the vacuum contractions give
\(\sum_{m\ge 0}\binom{2m}{m}(q_L^*q_R/4)^m=d^{-1/2}\).
The coherent-state completeness route in the original notes gives the
same result and is recorded in \cref{app:squeeze}.

The same derivatives of the bare kernel give the two operator insertions.
\begin{align}
 \frac{\mel{G_L}{b+b^\dagger}{G_R}}{h_{LR}} &= u+v,\\
 \frac{\mel{G_L}{b^\dagger b}{G_R}}{h_{LR}}
    &=\frac{\partial_x\partial_y\mathcal{F}}{\mathcal{F}}-1
      =uv+\frac{q_L^*q_R}{d}.
 \label{eq:transition-moments}
\end{align}
The subtraction comes from \(b^\dagger b=bb^\dagger-1\).
The prefactors \(c_L^*,c_R\) cancel in these ratios, which is why the
\(x,y\) derivatives act only on \(\mathcal{F}\).
At coincident packets, \(h_{LL}=1\), \(u=\alpha\), \(v=\alpha^*\), and
\(\langle b^\dagger b\rangle=|\alpha|^2+\sinh^2 r\).
Setting \(r=0\) recovers \cref{eq:coh-overlap,eq:coh-monomial} for these
operators.

\subsection{Local energy after contraction}
\label{sec:energy}

For the Hubbard-Holstein Hamiltonian
\begin{align}
  \hat H &= -t \sum_{\langle i,j \rangle, \sigma} (\hat c_{i \sigma}^\dagger \hat c_{i \sigma} + H.c.) + U \sum_i \hat n_{i\downarrow}\hat n_{i \uparrow} + \omega \sum_i \hat b_i^\dagger \hat b_i + g \sum_i \hat n_i (\hat b_i^\dagger + \hat b_i)  \label{eq:hubbard_holstein_hamiltonian} \\
  &= \hat H_t + \hat H_U + \hat H_\omega + \hat H_g, \notag
\end{align} 
let \(T_{ss'}=\mel{s}{\hat H_t}{s'}\). Write \(d_i(s)=n_{i\uparrow}(s)n_{i\downarrow}(s)\) and
\(n_i(s)=n_{i\uparrow}(s)+n_{i\downarrow}(s)\). As an operator on phonons,
the block \(H_{ss'}=\mel{s}{H}{s'}\) is
\begin{equation}
 H_{ss'}=T_{ss'}\id+\delta_{ss'}
 \sum_i\left[U d_i(s)\id+\omega_0 \hat b_i^\dagger \hat b_i+
               g n_i(s)(\hat b_i+\hat b_i^\dagger)\right].
 \label{eq:H-block}
\end{equation}
Inserting the fermion resolution of identity once gives
\begin{align}
 E &= \E_p[E_{\mathrm{loc}}(s)], \\
 E_{\mathrm{loc}}(s)
   &=\sum_{s'} \frac{A(s')}{A(s)}\, C_{H_{ss'}}(s,s') \notag\\
   &= \sum_{s'} T_{ss'}\,\frac{A(s')}{A(s)}\, C(s,s')
      + \sum_i\!\left[
        U d_i(s)+\omega_0 C_{b_i^\dagger b_i}(s,s)
        +g n_i(s) C_{b_i+b_i^\dagger}(s,s)\right].
 \label{eq:energy}
\end{align}
The first sum runs over hopping-connected configurations and includes
the fermionic sign through \(T_{ss'}\). A complex local energy can occur
for individual configurations. Its expectation is real.

In a joint fermion--phonon basis, this is the conditional average of
the ordinary local energy over the phonon configuration given \(s\).
The law of total variance reduces the single-draw variance by the
average conditional variance. Monte Carlo error also depends on
the autocorrelation of the remaining fermion sampler.

For the unnormalized conditional ansatz this is equivalently
\(\sum_{s'}[f(s')/f(s)]K_{H_{ss'}}(s,s')/Z_s\).
This form shows why the factor \(\sqrt{Z_{s'}}\) from the sampling
amplitude cancels that in the normalized kernel. The same construction
applies to any observable with known blocks \(B_{ss'}=\mel{s}{B}{s'}\).
A diagonal electronic observable such as \(n_i n_j\) is evaluated on
\(s\) directly. A hopping observable uses \(C(s,s')\). A phonon number
uses \(C_{b_i^\dagger b_i}(s,s)\).

\paragraph{One coherent packet per configuration.}
With its scalar weight absorbed into \(A(s)\), the phonon contribution
at fixed \(s\) is
\(\sum_i [\omega_0|\alpha_i(s)|^2+2g n_i(s)\Rea\alpha_i(s)]\).
Completing the square gives
\begin{equation}
 \omega_0|\alpha_i|^2+2g n_i\Rea\alpha_i
 =\omega_0\left|\alpha_i+\frac{g}{\omega_0}n_i\right|^2
     -\frac{g^2}{\omega_0}n_i^2 .
 \label{eq:atomic}
\end{equation}
For \(t=0\), the minimizing displacement is
\(\alpha_i(s)=-g n_i(s)/\omega_0\), the exact oscillator ground state in
that fermionic sector. Keeping this displacement at finite \(t\), a
single hop changes the occupations on two sites by one, so its real
phonon overlap is \(e^{-(g/\omega_0)^2}\).
This is the elementary Lang--Firsov overlap in the hopping channel.

\subsection{Energy gradient}
\label{sec:grad}

Use real variational coordinates \(\theta_\mu\), including real and
imaginary parts if the network outputs complex amplitudes. For
\(E=\langle\Psi|H|\Psi\rangle/\mathcal{N}\) and a
parameter-independent \(H\), the quotient rule gives
\begin{equation}
 \partial_\mu E=2\Rea\,
 \frac{\mel{\partial_\mu\Psi}{H-E}{\Psi}}{\mathcal{N}}.
 \label{eq:gradient}
\end{equation}
The following form applies automatic differentiation to the same
kernels used for the energy. Let \(\theta_0\) be the current parameters,
hold the ket fixed at \(\theta_0\), and let \(L\) be an independent bra
parameter vector. Define
\begin{align*}
 r_L(s) &=
       \frac{A_L(s)^*}{A_0(s)^*}, &
 C_B^{L,0}(s,s') &= \mel{\chi_s(L)}{B}{\chi_{s'}(\theta_0)},\\
 e^{L,0}(s)&=
 \sum_{s'}\frac{A_0(s')}{A_0(s)}\,C_{H_{ss'}}^{L,0}(s,s').
\end{align*}
Then the whole gradient is
\begin{equation}
 \partial_\mu E\big|_{\theta_0}
 =\left.\partial_{L_\mu}
 \left(2\Rea\,\E_{p_0}\!\left[
     r_L(s)\big(e^{L,0}(s)-E_0 C_{\id}^{L,0}(s,s)\big)
 \right]\right)\right|_{L=\theta_0}.
 \label{eq:surrogate}
\end{equation}
Indeed, multiplying the first term by
\(p_0(s)=|A_0(s)|^2/\mathcal{N}_0\) cancels both reference amplitudes
in the denominator, giving the bra-varied energy numerator divided by
\(\mathcal N_0\). The second term performs the same cancellation for
the norm. For a sampled batch, hold the configurations, the ket
parameters, the reference denominators and the energy estimate fixed
during this derivative. A same-configuration overlap equals one
after identifying the bra and ket parameters, so it must retain
its bra dependence in \cref{eq:surrogate}.

An equivalent form makes the force contributions explicit. Put
\begin{align}
 O_\mu(s)&=\partial_\mu\log A(s),&
 a_\mu(s)&=\braket{\chi_s}{\partial_\mu\chi_s},&
 \ell_\mu(s)&=O_\mu(s)+a_\mu(s), \notag\\
 \ket{D_\mu(s)}&=(\id-\ket{\chi_s}\bra{\chi_s})\ket{\partial_\mu\chi_s},&
 \ket{\eta_s}&=\sum_{s'}\frac{A(s')}{A(s)}H_{ss'}\ket{\chi_{s'}}.
 \label{eq:tangents}
\end{align}
Here the logarithmic derivative means \((\partial_\mu A)/A\), and
\(E_{\mathrm{loc}}(s)=\braket{\chi_s}{\eta_s}\).
The derivative of the amplitude-weighted conditional state is
\(A(s)[\ell_\mu(s)\ket{\chi_s}+\ket{D_\mu(s)}]\). Inserting it into
\cref{eq:gradient} gives
\begin{equation}
 \partial_\mu E
   = 2\Rea\,\E_p\!\left[
        \ell_\mu(s)^*(E_{\mathrm{loc}}(s)-E)
        +\braket{D_\mu(s)}{\eta_s}
      \right].
 \label{eq:force-decomp}
\end{equation}
The last term comes from changing the phonon state in a direction
orthogonal to itself. For the unnormalized conditional ansatz,
\(O_\mu=\partial_\mu\log f+\frac12\partial_\mu\log Z_s\).
The sampled gradient and metric formulas assume that the tangent has
support on the sampled configurations. If a configuration has exactly
zero amplitude but a nonzero derivative of the full component, its
contribution must be taken directly in the full Hilbert-space
expressions.

\subsection{Quantum geometric tensor and stochastic reconfiguration}
\label{sec:qgt}

The quantum geometric tensor (QGT) is
\begin{equation}
 S_{\mu\nu}=
 \frac{\braket{\partial_\mu\Psi}{\partial_\nu\Psi}}{\mathcal N}
 -\frac{\braket{\partial_\mu\Psi}{\Psi}\braket{\Psi}{\partial_\nu\Psi}}
        {\mathcal N^2}.
 \label{eq:qgt-global}
\end{equation}
In each electronic sector, the two pieces of the tangent in
\cref{eq:tangents} are orthogonal. Their inner products give the
uncentered \(\ell_\mu^*\ell_\nu\) term and the conditional QGT \(Q\),
while projection along the full state centers only \(\ell\). Thus
\begin{equation}
 S_{\mu\nu}
   = \E_p[(\ell_\mu-\E_p[\ell_\mu])^*(\ell_\nu-\E_p[\ell_\nu])]
     +\E_p[Q_{\mu\nu}(s)],
 \qquad Q_{\mu\nu}(s)=\braket{D_\mu(s)}{D_\nu(s)} .
 \label{eq:qgt-decomp}
\end{equation}
Expanding \(\ell=O+a\) gives the amplitude, bosonic phase-connection and
mixed covariances in the original notes. For a smooth locally real gauge
\(a_\mu=0\). \(Q_{\mu\nu}(s)\) can still be nonzero. For example, for a
product of single coherent states, differentiating
\cref{eq:coh-overlap} gives
\begin{align}
 a_\mu(s)&=\tfrac12\sum_i
     \left[\alpha_i^*\partial_\mu\alpha_i
            -\alpha_i\partial_\mu\alpha_i^*\right], &
 Q_{\mu\nu}(s)&=\sum_i
     (\partial_\mu\alpha_i)^*(\partial_\nu\alpha_i).
 \label{eq:qgt-coherent}
\end{align}
Under \(\ket{\chi_s}\to e^{\ii\varphi_s}\ket{\chi_s}\) and
\(A(s)\to e^{-\ii\varphi_s}A(s)\), the shifts in \(O_\mu\) and \(a_\mu\)
cancel and \(Q\) is invariant.

For general packets, duplicate the parameters on both sides of the
identity kernel. Keep \(L,R\) independent until both derivatives have
been taken.
\begin{equation}
 a_\nu(s)=
   \left.\partial_{R_\nu} C_{\id}^{L,R}(s,s)\right|_{L=R},\qquad
 Q_{\mu\nu}(s)=
   \left.\partial_{L_\mu}\partial_{R_\nu}
        \log K_{\id}^{L,R}(s,s)\right|_{L=R}.
 \label{eq:qgt-kernel}
\end{equation}
The mixed derivative is the mixed derivative of \(K\) divided by \(K\)
minus the product of its first derivatives divided by \(K^2\).
Normalization contributes only separate left and right terms to
\(\log C\), so the same mixed derivative can be taken on \(\log C\).
The logarithm is local around the coincident state, where \(K=Z_s>0\).

\subsubsection{Schur complement for a small phonon parameter block}
\label{sec:sr}

Suppose the variational coordinates split into disjoint blocks
\(\theta_f,\theta_b\) and \(\chi_s\) depends only on \(\theta_b\).
The amplitude \(A(s)\) may depend on both blocks, as for
\(A=f\sqrt{Z_s}\). For \(M\) sampled configurations let \(D_f,D_b\)
be the real matrices formed by taking the centered columns of \(\ell\)
in each block, dividing by \(\sqrt M\), and stacking the real and
imaginary parts vertically. Then \(D_x^{\mathsf T}D_y\) represents the
real part of the corresponding covariance block. Write
\(\overline Q_b=\Rea\,\widehat{\E}_p[Q_{bb}]\). With shift
\(\lambda>0\), the real SR equation for a direction \(\delta\) is
\begin{equation}
 \begin{pmatrix}
      \lambda I+D_f^{\mathsf T}D_f & D_f^{\mathsf T}D_b\\
      D_b^{\mathsf T}D_f & \lambda I+D_b^{\mathsf T}D_b+\overline Q_b
 \end{pmatrix}
 \begin{pmatrix}\delta_f\\\delta_b\end{pmatrix}
 =\begin{pmatrix}F_f\\F_b\end{pmatrix}.
 \label{eq:sr-block}
\end{equation}
For a descent direction one may take \(F=-\nabla E\).
If \(P=(\lambda I+D_f D_f^{\mathsf T})^{-1}\), the identity
\((\lambda I+D_f^{\mathsf T}D_f)^{-1}D_f^{\mathsf T}=D_f^{\mathsf T}P\)
gives the Schur complement and the reduced force.
\begin{align}
 (\lambda I+\overline Q_b+\lambda D_b^{\mathsf T}P D_b)\delta_b
       &=F_b-D_b^{\mathsf T}P D_f F_f,\label{eq:schur}\\
 \delta_f&=(\lambda I+D_f^{\mathsf T}D_f)^{-1}F_f
                 -D_f^{\mathsf T}P D_b\,\delta_b .
 \label{eq:schur-back}
\end{align}
In the first line the subtraction combines with the bosonic covariance
by \(I-D_f D_f^{\mathsf T}P=\lambda P\).
The first term in \cref{eq:schur-back} can be applied as
\(\lambda^{-1}(F_f-D_f^{\mathsf T}P D_f F_f)\). When
\(F_f=D_f^{\mathsf T}e\) for a sample-space force \(e\), as for the
covariance force on a purely fermionic block, it reduces to
\(D_f^{\mathsf T}Pe\). The inverse symbols denote linear solves.
The \(D_fD_f^{\mathsf T}\) solve is the MinSR step
\cite{chen2024}, while \(\overline Q_b\) contributes to the small
parameter-space solve. If the phonon network shares parameters with
the fermion network, \(Q\) can also occupy the other blocks and the
specialized reduction above then requires a different block solve.
For the normalized ansatz with \(A=f(\theta_f)\) and real conditional
phonon states, \(\ell_b=a_b=0\) and the two blocks decouple.

\subsection{Fermion--phonon second R\'enyi entropy}
\label{sec:renyi}

For the normalized density matrix
\(\rho=\ket{\Psi}\bra{\Psi}/\mathcal N\), tracing out the fermions gives
a statistical mixture of the conditional states.
\begin{equation}
 \rho_b=\sum_s p(s)\ket{\chi_s}\bra{\chi_s},\qquad
 (\rho_f)_{ss'}=
       \frac{A(s)A(s')^*}{\mathcal N}\,C(s',s).
 \label{eq:reduced}
\end{equation}
Both reduced density matrices have the same nonzero spectrum.
Squaring \(\rho_b\) and taking the trace gives
\begin{equation}
 \mathcal P=\Tr(\rho_b^2)=
 \E_{s,s'\,\mathrm{ind.}\sim p}|C(s,s')|^2,\qquad
 S_2=-\log\mathcal P.
 \label{eq:purity}
\end{equation}
Equivalently the pair observable is
\(\lvert K_{\id}(s,s')\rvert^2/(Z_s Z_{s'})\).
This is the two-copy swap identity for purity with the phonon
contraction already done \cite{hastings2010}.
The pair observable is the pure-state fidelity of two phonon clouds
and lies in \([0,1]\). When all clouds on the support have the same
ray, the fidelity is one and \(S_2=0\). For a single coherent packet,
\[
 |C(s,s')|^2=\exp\!\left[-\sum_i|\alpha_i(s)-\alpha_i(s')|^2\right].
\]

With two independent chains at stationarity, averaging one matched
pair per index gives the purity estimator with \(M\) overlap evaluations.
\begin{equation}
 \widehat{\mathcal P}
   =\frac1M\sum_{j=1}^{M}|C(s_j,s'_j)|^2 .
 \label{eq:purity-mc}
\end{equation}
All \(M^2\) cross-chain pairs can also be averaged. Each overlap costs
\(O(K^2N)\) to contract once the network outputs are available.
For \(M\) independent draws from a single sample set, the average over
the distinct index pairs \(i\ne j\), normalized by \(M(M-1)\), is
unbiased for \(\mathcal P\).
The all-pairs estimator in the notes includes \(M\) self-pairs with
fidelity one, giving an upward purity bias of
\((1-\mathcal P)/M\) even for independent draws.
In a single Markov chain, distinct indices can still be correlated.
Taking the logarithm of the estimated purity also introduces
finite-sample bias.

\subsection{Relation to prior work}
\label{sec:references}

Coherent-state overlap and normal-ordering algebra are standard
\cite{glauber1963}. The exponential-quadratic kernel belongs to the
more general Bargmann treatment of Gaussian states
\cite{john2021}. Cataudella \emph{et al.} used coherent-state expansions
for optimized phonon exact diagonalization \cite{cataudella2004}.
Bera \emph{et al.} used conditional multimode coherent superpositions
and analytical contractions in the spin--boson model
\cite{bera2014}.
Walsh, Boettcher and Marsiglio give a close example of the analytic
algebra for the one-electron Holstein polaron. Their Appendix A uses
coherent overlaps and annihilation eigenvalues to construct the
overlap and Hamiltonian matrices \cite{walsh2026}.
For a direct many-electron VMC precedent, Ohgoe and Imada sample both
the electronic configuration and phonon coordinates
\cite{ohgoe2014}. Mahajan \emph{et al.} develop neural VMC for
electron--phonon problems in a joint configuration basis
\cite{mahajan2025}. Here the conditional packet representation
contracts the phonons analytically before sampling fermions.

\clearpage
\appendix

\subsection{Disentangling and the Gaussian integral}
\label[appendix]{app:squeeze}

Set \(K_+=(b^\dagger)^2/2\), \(K_-=b^2/2\), and
\(K_0=(b^\dagger b+1/2)/2\). The commutators
\([K_+,K_-]=-2K_0\) and \([K_0,K_\pm]=\pm K_\pm\)
allow the two-by-two representation
\[
 K_+\mapsto\begin{pmatrix}0&1\\0&0\end{pmatrix},\qquad
 K_-\mapsto\begin{pmatrix}0&0\\-1&0\end{pmatrix},\qquad
 K_0\mapsto\tfrac12\begin{pmatrix}1&0\\0&-1\end{pmatrix}.
\]
Exponentiating \(\zeta^* K_--\zeta K_+\) in this representation gives
diagonal entries \(\cosh r\) and off-diagonal entries
\(-e^{\pm\ii\vartheta}\sinh r\). Multiplying the upper triangular,
diagonal and lower triangular factors gives the operator identity
\begin{equation}
 S(\zeta)=\exp(qK_+)\exp[-2\log(\cosh r)K_0]\exp(-q^*K_-).
 \label{eq:disentangle}
\end{equation}
The right factor acts trivially on the vacuum and the middle factor
gives \((\cosh r)^{-1/2}\). Displacing the result uses
\(D(\alpha)b^\dagger D(\alpha)^\dagger=b^\dagger-\alpha^*\) and
\(\ket{\alpha}=e^{-|\alpha|^2/2}e^{\alpha b^\dagger}\ket0\), so the
quadratic in \(b^\dagger-\alpha^*\) gives precisely the linear term and
prefactor in \cref{eq:gaussian-state}.

Alternatively, insert the coherent-state resolution of the identity
\(\int_{\mathbb C}\frac{\dd^2 z}{\pi}\ket z\bra z=\id\)
between the bare exponentials. The integral is
\begin{equation}
 \mathcal F(x,y)=
 \int_{\mathbb C}\frac{\dd^2 z}{\pi}\exp\!\left[
    -|z|^2+\tfrac12 q_L^*z^2+xz+
       \tfrac12 q_R(z^*)^2+yz^*\right].
 \label{eq:integral}
\end{equation}
Completing the square in \(\Rea z,\Ima z\) gives the determinant
\(d=1-q_L^*q_R\) and the exponent in \cref{eq:gauss-overlap}.

\subsection{Single-mode cavity Peierls coupling}
\label[appendix]{app:cavity}

For a link with the photon operator
\(e^{\ii\kappa(a+a^\dagger)}=D(\ii\kappa)\), \(\kappa\in\mathbb R\),
the block formula in \cref{eq:energy} uses the displaced overlap
on that link. For coherent packets, normal ordering gives
\begin{equation}
 \mel{\alpha_L}{e^{\ii\kappa(a+a^\dagger)}}{\alpha_R}
   = h_{LR}\exp\!\left[-\kappa^2/2+
                         \ii\kappa(\alpha_L^*+\alpha_R)\right].
 \label{eq:peierls}
\end{equation}
The reverse hopping uses \(-\kappa\).
For squeezed packets the Weyl relation gives, with \(\xi=\ii\kappa\),
\[
 D(\xi)\ket{G_R}
   = e^{(\xi\alpha_R^*-\xi^*\alpha_R)/2}
        D(\alpha_R+\xi) S(\zeta_R)\ket{0}.
\]
Hence evaluate \cref{eq:gauss-overlap} with
\(\alpha_R\to\alpha_R+\xi\), then multiply by the indicated phase.
For a packet superposition the contraction is the same
\(W(s)^\dagger G_B(s,s')W(s')\) as above.
This analytically contracts the entire oscillator Hilbert space for
the chosen ansatz. A finite Fock-basis sum contracts the projected
state and can also be used for arbitrary photon-dependent amplitudes.

\clearpage
\subsection*{Changes from the supplied notes}
\begin{itemize}
 \item Equations (6)--(8) and (30) use the same symbol for unitary
 squeezing and the quadratic coefficient, omit the displacement
 shift of the linear coefficient, and show a positive sign in the overlap
 prefactor exponent. The corrected kernel has unit self-overlap
 and retains the coherent-state phase.
 \item Equations (22), (26)--(29), and (32)--(33) are covered by the
 corrected normal-ordering sign, the explicit Hubbard prefactor \(U\),
 a single Gram contraction, and the energy gradient including its
 norm derivative. The sampling measure follows the state convention.
 \item Equations (39)--(41) give the starting block elimination. The
 reduced expressions here retain the boson block and force and subtract
 the corresponding fermionic correction. Sample centering and
 regularization are explicit.
 \item Equations (82)--(83) now allow complex weights and displacements
 via the full kernel. Equation (93) is replaced by independent pair
 sampling or the distinct-index estimator for independent draws.
\end{itemize}

\begin{thebibliography}{9}

\bibitem{glauber1963}
R.~J. Glauber, \emph{Coherent and Incoherent States of the Radiation Field},
Phys.\ Rev.\ \textbf{131}, 2766 (1963).
\href{https://doi.org/10.1103/PhysRev.131.2766}{DOI}.

\bibitem{chen2024}
A. Chen and M. Heyl, \emph{Empowering deep neural quantum states through
efficient optimization}, Nat.\ Phys.\ \textbf{20}, 1476--1481 (2024).
\href{https://doi.org/10.1038/s41567-024-02566-1}{DOI}.

\bibitem{hastings2010}
M.~B. Hastings, I. Gonz\'alez, A.~B. Kallin, and R.~G. Melko,
\emph{Measuring Renyi entanglement entropy in quantum Monte Carlo
simulations}, Phys.\ Rev.\ Lett.\ \textbf{104}, 157201 (2010).
\href{https://doi.org/10.1103/PhysRevLett.104.157201}{DOI}.

\bibitem{john2021}
T.~C. John and K.~R. Parthasarathy,
\emph{A common parametrization for finite mode Gaussian states, their
symmetries, and associated contractions with some applications},
J.\ Math.\ Phys.\ \textbf{62}, 022102 (2021).
\href{https://arxiv.org/abs/1911.06555}{arXiv:1911.06555}.

\bibitem{cataudella2004}
V. Cataudella, G. De Filippis, F. Martone, and C.~A. Perroni,
\emph{Variational approach to the optimized phonon technique for
electron-phonon problems}, Phys.\ Rev.\ B \textbf{70}, 193105 (2004).
\href{https://doi.org/10.1103/PhysRevB.70.193105}{DOI}.

\bibitem{bera2014}
S. Bera, A. Nazir, A.~W. Chin, H.~U. Baranger, and S. Florens,
\emph{Generalized multipolaron expansion for the spin-boson model:
Environmental entanglement and the biased two-state system},
Phys.\ Rev.\ B \textbf{90}, 075110 (2014).
\href{https://doi.org/10.1103/PhysRevB.90.075110}{DOI}.

\bibitem{walsh2026}
C.~M. Walsh, I. Boettcher, and F. Marsiglio,
\emph{Coherent-state ansatz for the Holstein polaron in one and two
dimensions} (2026).
\href{https://arxiv.org/html/2603.07290v1}{arXiv:2603.07290v1}.

\bibitem{ohgoe2014}
T. Ohgoe and M. Imada, \emph{Variational Monte Carlo method for
electron-phonon coupled systems}, Phys.\ Rev.\ B \textbf{89}, 195139 (2014).
\href{https://doi.org/10.1103/PhysRevB.89.195139}{DOI}.

\bibitem{mahajan2025}
A. Mahajan, P.~J. Robinson, J. Lee, and D.~R. Reichman,
\emph{Structure and dynamics of electron-phonon coupled systems using
neural quantum states}, Phys.\ Rev.\ B \textbf{112}, 184310 (2025).
\href{https://doi.org/10.1103/p8qz-p45d}{DOI}.

\end{thebibliography}
\end{document}


